\documentclass[11pt,a4paper]{article}

% ─── Encodage & langue ────────────────────────────────────────────────────────
\usepackage[utf8]{inputenc}
\usepackage[T1]{fontenc}
\usepackage[french]{babel}

% ─── Mise en page ─────────────────────────────────────────────────────────────
\usepackage[top=2cm, bottom=2cm, left=2.2cm, right=2.2cm]{geometry}
\usepackage{parskip}
\setlength{\parindent}{0pt}

% ─── Couleurs & boîtes ────────────────────────────────────────────────────────
\usepackage[dvipsnames,table]{xcolor}
\definecolor{springgreen}{RGB}{109,179,63}
\definecolor{j2eered}{RGB}{190,45,45}
\definecolor{codebg}{RGB}{248,248,248}
\definecolor{sectionblue}{RGB}{0,70,150}
\definecolor{headerbg}{RGB}{30,30,30}

% ─── Tableaux ─────────────────────────────────────────────────────────────────
\usepackage{booktabs}
\usepackage{tabularx}
\usepackage{array}
\newcolumntype{L}[1]{>{\raggedright\arraybackslash}p{#1}}

% ─── Listings ─────────────────────────────────────────────────────────────────
\usepackage{listings}
\lstset{
  basicstyle=\ttfamily\footnotesize,
  backgroundcolor=\color{codebg},
  frame=single, framesep=3pt,
  rulecolor=\color{gray!40},
  breaklines=true, keepspaces=true,
  columns=fullflexible,
  showstringspaces=false,
  keywordstyle=\color{sectionblue}\bfseries,
  commentstyle=\color{gray}\itshape,
  stringstyle=\color{BrickRed},
  numbers=none
}

% ─── Titres ───────────────────────────────────────────────────────────────────
\usepackage{titlesec}
\titleformat{\section}{\large\bfseries\color{sectionblue}}{\thesection}{0.8em}{}[\vspace{-4pt}\color{sectionblue}\rule{\linewidth}{0.6pt}]
\titleformat{\subsection}{\normalsize\bfseries\color{j2eered!80!black}}{\thesubsection}{0.8em}{}
\titlespacing*{\section}{0pt}{14pt}{4pt}
\titlespacing*{\subsection}{0pt}{8pt}{2pt}

% ─── Divers ───────────────────────────────────────────────────────────────────
\usepackage{enumitem}
\usepackage{hyperref}
\hypersetup{colorlinks=true, linkcolor=sectionblue, pdftitle={Architecture Game Portal}}

% ─── Commandes utilitaires ────────────────────────────────────────────────────
\newcommand{\spring}[1]{\textcolor{springgreen}{\textbf{[Spring]}} #1}
\newcommand{\jee}[1]{\textcolor{j2eered}{\textbf{[J2EE]}} #1}
\newcommand{\code}[1]{\texttt{\small #1}}

% ─── En-tête de tableau coloré ────────────────────────────────────────────────
\newcommand{\theadrow}[2]{%
  \rowcolor{headerbg}\textcolor{white}{\textbf{#1}} & \textcolor{white}{\textbf{#2}}\\}
\newcommand{\theadrowT}[3]{%
  \rowcolor{headerbg}\textcolor{white}{\textbf{#1}} & \textcolor{white}{\textbf{#2}} & \textcolor{white}{\textbf{#3}}\\}

% ══════════════════════════════════════════════════════════════════════════════
\begin{document}
% ══════════════════════════════════════════════════════════════════════════════

% ─── Page de titre ────────────────────────────────────────────────────────────
\begin{center}
  {\LARGE\bfseries Architecture du projet Game Portal}\\[6pt]
  {\large\color{sectionblue} Spring Boot 3.4.5 vs Servlet / J2EE}\\[4pt]
  {\small Documentation pédagogique — 2026}
\end{center}
\vspace{6pt}
\hrule height 1.5pt
\vspace{10pt}

\tableofcontents
\vspace{6pt}
\hrule
\newpage

% ══════════════════════════════════════════════════════════════════════════════
\section{Vue d'ensemble}
% ══════════════════════════════════════════════════════════════════════════════

Le projet est une application full-stack : \textbf{Angular :4200} (frontend) + \textbf{Spring Boot :8080} (backend) + \textbf{MySQL} (base \code{game\_portal}).

\begin{lstlisting}[language={}]
Angular :4200  -->  SecurityFilterChain (Spring Security)
                        |
               +---------+----------+
               |         |          |
         AuthController  GameController  ScoreController  (Spring MVC)
               |         |          |
         AuthService   GameIntegrationService             (Services)
               |         |          |
         UserRepo    GameRepo    ScoreRepo                (Spring Data JPA)
               |
           MySQL (tables : users, games, scores)
\end{lstlisting}

\begin{table}[h!]
\centering
\renewcommand{\arraystretch}{1.35}
\begin{tabularx}{\linewidth}{L{4.8cm} X}
\theadrow{Dépendance Spring Boot}{Rôle}
\code{spring-boot-starter-web}       & Spring MVC — remplace le mapping Servlet dans \code{web.xml} \\
\code{spring-boot-starter-data-jpa}  & Hibernate/JPA — remplace JDBC + \code{persistence.xml} \\
\code{spring-boot-starter-security}  & Filtres, BCrypt, contrôle d'accès — remplace les filtres manuels \\
\code{spring-boot-starter-validation}& \code{@Valid}, \code{@NotBlank} — remplace les validations \code{if/else} \\
\code{jjwt-api / jjwt-impl}          & Génération et validation des tokens JWT \\
\code{lombok}                        & Génère getters/setters/constructeurs automatiquement \\
\bottomrule
\end{tabularx}
\caption{Dépendances principales}
\end{table}

% ══════════════════════════════════════════════════════════════════════════════
\section{Spring MVC — Couche Contrôleurs}
% ══════════════════════════════════════════════════════════════════════════════

\spring{Un seul \code{@RestController} regroupe tous les endpoints d'un domaine.}

\jee{Chaque endpoint = une classe \code{HttpServlet} + une entrée dans \code{web.xml}.}

\begin{table}[h!]
\centering
\renewcommand{\arraystretch}{1.35}
\begin{tabularx}{\linewidth}{L{4.8cm} X}
\theadrow{Annotation Spring MVC}{Équivalent J2EE}
\code{@RestController}               & \code{extends HttpServlet} + \code{resp.setContentType("application/json")} \\
\code{@RequestMapping("/api/auth")}  & \code{<url-pattern>} dans \code{web.xml} \\
\code{@PostMapping}, \code{@GetMapping} & \code{doPost()}, \code{doGet()} dans la Servlet \\
\code{@RequestBody RegisterRequest}  & Lecture manuelle via \code{request.getReader()} + \code{new JSONObject()} \\
\code{@Valid}                        & Chaîne \code{if/else} de validation répétée dans chaque Servlet \\
\code{@RequestParam String game}     & \code{request.getParameter("game")} \\
\code{@AuthenticationPrincipal User} & \code{(User) request.getAttribute("currentUser")} (propagé manuellement) \\
\code{ResponseEntity.ok(body)}       & \code{resp.setStatus(200); resp.getWriter().write(json)} \\
Jackson (auto)                       & \code{StringBuilder} JSON + méthode \code{escapeJson()} — fragile \\
\bottomrule
\end{tabularx}
\caption{Annotations Spring MVC vs J2EE}
\end{table}

% ══════════════════════════════════════════════════════════════════════════════
\section{Spring Data JPA — Couche Persistance}
% ══════════════════════════════════════════════════════════════════════════════

\spring{\code{JpaRepository} : interface de 5 lignes, Spring génère tout le SQL.}

\jee{\code{UserDAO} : classe de 120 lignes JDBC (connexion, \code{PreparedStatement}, \code{ResultSet}, fermeture).}

\begin{table}[h!]
\centering
\renewcommand{\arraystretch}{1.35}
\begin{tabularx}{\linewidth}{L{5cm} X}
\theadrow{Spring Data JPA}{Équivalent JDBC J2EE}
\code{findByUsername(String u)}            & \code{SELECT * FROM users WHERE username=?} + mapping \code{ResultSet} \\
\code{existsByEmail(String e)}             & \code{SELECT COUNT(*) FROM users WHERE email=?} \\
\code{save(user)}                          & \code{INSERT} ou \code{UPDATE} avec \code{PreparedStatement} \\
\code{findByUserOrderByPlayedAtDesc(user)} & \code{SELECT * FROM scores WHERE user\_id=? ORDER BY played\_at DESC} \\
Pool HikariCP (auto)                       & \code{DriverManager.getConnection()} — nouvelle connexion à chaque appel \\
\code{@Transactional}                      & \code{conn.commit()} / \code{conn.rollback()} manuels \\
\midrule
\code{@Entity @Table(name="users")}        & Script SQL \code{CREATE TABLE users (...)} exécuté manuellement \\
\code{@Column(unique=true, length=50)}     & \code{VARCHAR(50) UNIQUE} dans le DDL \\
\code{@GeneratedValue(IDENTITY)}           & \code{AUTO\_INCREMENT} + \code{getGeneratedKeys()} \\
\code{@ManyToOne @JoinColumn("user\_id")}  & \code{FOREIGN KEY} + \code{JOIN} SQL + mapping des deux objets \\
\code{@CreationTimestamp}                  & \code{NOW()} dans l'\code{INSERT} \\
\code{ddl-auto=update}                     & Script \code{schema.sql} exécuté manuellement \\
\bottomrule
\end{tabularx}
\caption{Spring Data JPA / Hibernate vs JDBC J2EE}
\end{table}

% ══════════════════════════════════════════════════════════════════════════════
\section{Spring Security — Couche Sécurité}
% ══════════════════════════════════════════════════════════════════════════════

\subsection{Flux d'authentification (JWT)}

\begin{lstlisting}[language={}]
POST /api/auth/login
  --> AuthenticationManager.authenticate()
        --> UserDetailsServiceImpl.loadUserByUsername()  [DB]
        --> BCryptPasswordEncoder.matches()
  --> jwtUtil.generateToken(user)  -->  JWT signe HS256 (24h)
  <-- { token, username, role }

GET /api/scores/history  + Header: Authorization: Bearer <token>
  --> JwtAuthFilter
        --> jwtUtil.extractUsername(token)
        --> jwtUtil.validateToken(token, userDetails)
        --> SecurityContextHolder.setAuthentication(authToken)
  --> @AuthenticationPrincipal User currentUser  (injecte depuis SecurityContext)
\end{lstlisting}

\subsection{Tableau comparatif}

\begin{table}[h!]
\centering
\renewcommand{\arraystretch}{1.35}
\begin{tabularx}{\linewidth}{L{5cm} X}
\theadrow{Spring Security}{Équivalent J2EE}
\code{SecurityFilterChain} (1 classe) & \code{web.xml} + 3--4 classes \code{Filter} séparées \\
\code{.permitAll()}                    & Routes exclues du \code{<url-pattern>} dans \code{web.xml} \\
\code{.hasRole("ADMIN")}               & \code{if (!"ADMIN".equals(session.getAttribute("role")))} répété partout \\
\code{.authenticated()}                & \code{if (session == null) sendError(401)} dans chaque Servlet \\
\code{SessionCreationPolicy.STATELESS} & Pas d'équivalent — session HTTP créée par défaut \\
\code{BCryptPasswordEncoder}           & \code{MessageDigest SHA-256} sans sel — vulnérable aux rainbow tables \\
\code{OncePerRequestFilter}            & \code{javax.servlet.Filter} sans garantie d'exécution unique \\
\code{SecurityContextHolder}           & \code{request.setAttribute("currentUser", user)} manuel \\
\code{@PreAuthorize("hasRole('ADMIN')")} & Vérification \code{if/else} répétée dans chaque Servlet admin \\
\bottomrule
\end{tabularx}
\caption{Spring Security vs filtres J2EE}
\end{table}

\subsection{JWT vs Session HTTP}

\begin{table}[h!]
\centering
\renewcommand{\arraystretch}{1.35}
\begin{tabularx}{\linewidth}{L{3.5cm} X X}
\theadrowT{Critère}{JWT (Spring)}{Session HTTP (J2EE)}
Identifiant        & En-tête \code{Authorization: Bearer} & Cookie \code{JSESSIONID} \\
Stockage serveur   & Aucun (stateless)             & En mémoire / Redis \\
Scalabilité        & Horizontale sans config       & Sticky sessions ou Redis \\
CSRF               & Impossible (pas de cookie)    & Possible \\
Révocation         & Difficile (blacklist)         & Immédiate (\code{session.invalidate()}) \\
Contenu            & Username + rôle inclus        & ID de session seulement \\
\bottomrule
\end{tabularx}
\caption{JWT vs Session HTTP}
\end{table}

% ══════════════════════════════════════════════════════════════════════════════
\section{Services — Logique Métier}
% ══════════════════════════════════════════════════════════════════════════════

\begin{table}[h!]
\centering
\renewcommand{\arraystretch}{1.35}
\begin{tabularx}{\linewidth}{L{5cm} X}
\theadrow{Spring (AuthService)}{J2EE (RegisterServlet)}
\code{@Transactional} — rollback auto    & Pas de transaction — état incohérent possible \\
\code{passwordEncoder.encode(pwd)}       & \code{MessageDigest SHA-256} sans sel \\
\code{userRepository.save(user)}         & \code{INSERT} JDBC \textasciitilde15 lignes \\
\code{authenticationManager.authenticate()} & \code{SELECT} + comparaison hash manuelle \\
2 méthodes claires, \textasciitilde60 lignes & 2 Servlets \textasciitilde140 lignes au total \\
\midrule
\multicolumn{2}{l}{\textbf{GameIntegrationService — Appel API externe}} \\
\midrule
\code{restTemplate.getForObject(url, JsonNode.class)} — 1 ligne & \code{HttpURLConnection} + \code{BufferedReader} — \textasciitilde25 lignes \\
Jackson désérialise en \code{JsonNode} automatiquement & \code{new JSONArray(body)} — parsing champ par champ \\
\code{@Slf4j} : \code{log.info/error()} avec niveau et contexte & \code{System.out.println()} — sans niveau ni contexte \\
\bottomrule
\end{tabularx}
\caption{Spring Services vs logique dans les Servlets J2EE}
\end{table}

% ══════════════════════════════════════════════════════════════════════════════
\section{Bean Validation et Gestion des Erreurs}
% ══════════════════════════════════════════════════════════════════════════════

\spring{\code{@NotBlank}, \code{@Email}, \code{@Size} sur le DTO + \code{@Valid} dans le contrôleur.}

\jee{\textasciitilde20 lignes de \code{if/else} par Servlet, dupliquées partout.}

\begin{table}[h!]
\centering
\renewcommand{\arraystretch}{1.35}
\begin{tabularx}{\linewidth}{L{3.5cm} L{2.5cm} X}
\theadrowT{Exception Spring}{Statut HTTP}{Équivalent J2EE}
\code{MethodArgumentNotValidException} & 400 + détails & \code{if/else} de validation dans chaque Servlet \\
\code{AccessDeniedException}           & 403 JSON      & \code{sendError(403)} répété dans chaque Servlet admin \\
\code{AuthenticationException}         & 401 JSON      & \code{if (session==null) sendError(401)} \\
\code{Exception} (filet)               & 500 JSON      & Page HTML Tomcat avec stacktrace — fuite d'info \\
\bottomrule
\end{tabularx}
\caption{\code{@RestControllerAdvice} vs gestion d'erreurs J2EE}
\end{table}

% ══════════════════════════════════════════════════════════════════════════════
\section{Lombok}
% ══════════════════════════════════════════════════════════════════════════════

\begin{table}[h!]
\centering
\renewcommand{\arraystretch}{1.35}
\begin{tabularx}{\linewidth}{L{4.8cm} X}
\theadrow{Annotation Lombok}{Code Java généré}
\code{@Data}               & Getters, setters, \code{equals()}, \code{hashCode()}, \code{toString()} \\
\code{@Builder}            & Pattern Builder : \code{User.builder().username("a").build()} \\
\code{@RequiredArgsConstructor} & Constructeur pour champs \code{final} — injection de dépendances \\
\code{@Slf4j}              & \code{private static final Logger log = LoggerFactory.getLogger(…)} \\
\bottomrule
\end{tabularx}
\caption{Annotations Lombok}
\end{table}

% ══════════════════════════════════════════════════════════════════════════════
\section{Bilan général}
% ══════════════════════════════════════════════════════════════════════════════

\begin{table}[h!]
\centering
\renewcommand{\arraystretch}{1.4}
\begin{tabularx}{\linewidth}{L{5.5cm} L{2.8cm} X}
\theadrowT{Fonctionnalité}{Spring Boot}{J2EE équivalent}
Démarrage serveur              & 1 ligne \code{main()} & Déploiement \code{.war} + config Tomcat \\
Sécurité complète              & 60 lignes             & \code{web.xml} + 3 Filter + code inline — 400+ lignes \\
CRUD utilisateurs              & 10 lignes             & \code{UserDAO} — 120 lignes JDBC \\
1 endpoint REST                & 3 lignes              & 1 Servlet + \code{web.xml} + JSON manuel \\
Validation des entrées         & 3 annotations         & 20 \code{if/else} par Servlet \\
Gestion des erreurs            & 1 classe globale      & \code{try/catch} copié dans chaque Servlet \\
Appel API externe              & 1 ligne               & \code{HttpURLConnection} — 25 lignes \\
Encodage mot de passe          & BCrypt (sel auto)     & SHA-256 sans sel — vulnérable \\
\bottomrule
\end{tabularx}
\caption{Comparaison globale Spring Boot vs J2EE}
\end{table}

\vspace{8pt}
\noindent\textbf{Avantages clés de Spring Boot :}
\begin{enumerate}[leftmargin=*, itemsep=2pt]
  \item \textbf{Configuration centralisée} : sécurité, CORS et routes dans une seule classe.
  \item \textbf{IoC / Injection de dépendances} : classes découplées et testables unitairement.
  \item \textbf{Spring Data JPA} : SQL généré automatiquement, HikariCP intégré.
  \item \textbf{Sécurité robuste} : BCrypt résistant aux rainbow tables, JWT stateless scalable.
  \item \textbf{Séparation des responsabilités} : Controller $\to$ Service $\to$ Repository — pas de logique dans les Servlets.
  \item \textbf{Erreurs unifiées} : \code{@RestControllerAdvice} retourne toujours du JSON structuré.
\end{enumerate}

\end{document}
